\documentclass[withoutpreface,bwprint]{cumcmthesis} %去掉封面与编号页，电子版提交的时候使用。

\usepackage{tabularx}
\usepackage{needspace}
\usepackage{longtable}
\usepackage{cite}
\usepackage[framemethod=TikZ]{mdframed}
\usepackage{url}   % 网页链接
\usepackage{subcaption} % 子标题
\usepackage{titlesec}    % 控制标题样式
\usepackage{doi}
\usepackage{amsmath}
\numberwithin{equation}{section}

% 强制公式编号为阿拉伯数字
\renewcommand{\theequation}{\arabic{section}.\arabic{equation}}

% 自定义四级标题计数器
\newcounter{fourthlevel}[subsubsection]
\renewcommand{\thefourthlevel}{\thesubsubsection.\arabic{fourthlevel}}
\makeatletter
\@addtoreset{paragraph}{subsubsection}
\makeatother

\titleformat{\paragraph}[block]
  {\heiti\large\bfseries}
  {\thesubsubsection.\arabic{paragraph}}
  {1em}
  {}
\titlespacing{\paragraph}{0pt}{1ex plus 0.5ex minus .2ex}{1em}

% 让 paragraph 参与编号
\setcounter{secnumdepth}{4}

\title{基于联合场景优化的蔬菜类商品的定价与补货策略研究}



\tihao{C}
\baominghao{}
\schoolname{}
\membera{ }
\memberb{ }
\memberc{ }
\supervisor{ }
\yearinput{2025}
\monthinput{09}
\dayinput{04}

\begin{document}

 \maketitle
 \begin{abstract}
蔬菜类商品在生鲜商超的日常经营中，通常具有保质期短、需求不稳定的特点。若未能在当日内售出，往往会无法再售，导致商超面临损失。
但由于各类蔬菜品种、产地等因素存在较大差异，蔬菜销售情况受到多种因素的影响，商家需要在不知道可供应单品以及进货价格的条件下在凌晨对当天各蔬菜的补货数量进行决策。
因此，本文结合销量时间规律、季节性供应条件并依据市场需求，确定合理的蔬菜补货方案以及定价策略，减少商超损失。

\textbf{针对问题一，}采用\textbf{两部分布模型}分析销量分布，并通过 \textbf{AIC 与 KS 检验确定最优分布}；进一步利用月度销量占比和活跃率构建季节特征，通过\textbf{ K-means }将单品划分为五种季节销售模式。在此基础上，从季节同步性、销量共变、份额变化和销售活跃期四个方面分析相互关系。结果表明，\textbf{六品类总销量相关性较强（$[0.415,0.692]$），但份额相关性最低达 $-0.535$，表现出总量同步与结构竞争并存的特征，而单品间稳定相关较少。}

\textbf{针对问题二}在品类层面分析销售总量与成本加成定价的关系，并给出 2023 年 7 月 1—7 日的日补货总量与定价策略。采用\textbf{部分线性双重机器学习}（DML）估计六品类的条件价格响应,以第一问品类间正相关为依据，按共同日期拼接 12 维残差向量生成 600 个联合 7 日场景，保留品类间相关、需求—成本相关与周内自相关。以均值—下尾 CVaR 为\textbf{优化目标}，周级 SLSQP 联合求解 42 个补货量与 42 个加成率。\textbf{主方案一周补货 2048.33 kg，经营利润期望 4850.30 元，相对历史基线提高约 2179 元；}最差 10\% 经营利润均值由负转正，收益与尾部风险同时改善。

\textbf{针对问题三：}在问题二的品类级决策基础上，进一步制定单品补货与定价方案，
要求每日可售单品数控制在27--33个，且单品订货量不低于$2.5\,\mathrm{kg}$。
采用\textbf{经验贝叶斯动态收缩}估计品类内部单品需求份额，并结合问题二得到的
品类需求、价格弹性和成本信息构建单品需求场景。随后采用\textbf{两阶段词典序优化}，
优先提高六个品类的需求满足率，再在服务水平基本保持不变的条件下优化期望利润
与尾部收益。对$K=27,\ldots,33$分别求解并在完整600个场景下进行评价，最终选取
$K=33$，得到33个可售单品、总订货量$362.26\,\mathrm{kg}$。需求满足率为
$77.81\%$，\textbf{期望利润为$163.75$元，}较近期销量基线分别提高13.20个百分点和
107.27元，整体优于常规经验策略。

\textbf{针对问题四：}结合前述模型分析结果，提出商超数据采集建议。重点补充门店客流、天气与节假日以及竞争门店价格等信息，以减少缺货造成的需求偏差，改善价格弹性识别和需求预测，为后续补货与定价决策提供更准确的数据支持。


\keywords{两部分布；K-means 季节聚类；条件价格响应；联合残差场景；均值—下尾 CVaR；两阶段词典序优化}
\end{abstract}


%\tableofcontents

%\newpage

\section{问题重述}

生鲜商超中，蔬菜类商品具有保鲜时期短、品质下降快的特点。当日未售完的蔬菜，次日已不适宜继续销售。相关研究表明，水果和蔬菜是超市易损耗商品的重要组成部分，不同蔬菜单品的零售损耗率存在较大差异\textsuperscript{\cite{buzby2015produce}}。
因此，合理的库存与销售管理有助于兼顾商品新鲜度、销售量和损耗控制\textsuperscript{\cite{broekmeulen2019foodwaste}}。

商超蔬菜采用成本加定价的方法进行销售，同时对运损和品质下降的蔬菜进行打折销售。但由于蔬菜采购通常在凌晨进行，
商家在不明确当天具体供应单品及价格的情况下需要制定各类蔬菜的补货计划。因此，对市场需求进行准确预测分析是合理制定定价策略的重要基石。
从需求方面看，蔬菜的销售量与时间之间也存在一定联系；从供给角度看，每年4月至10月的蔬菜可供应品种较多。同时，在有限的商超空间内，合理的销售组合有重要意义。

已知附件1提供某商超经营的6个蔬菜品类及其商品信息，附件2与附件3分别给出了2020年7月1日至2023年6月30日期间各商品的销售流水数据和批发价格数据，附件4则提供了各商品近期的损耗率信息。
现结合所给数据以及实际经营情况建立数学模型解决以下问题：

问题1：不同蔬菜品类以及单品之间可能存在一定的联系。完成对蔬菜品类与单品销售量的分布规律以及相互关系的分析。

问题2：在商超以蔬菜品类为单位进行补货的情况下，分析不同品类销售总量与成本加成定价的关系，并制定2023年7月1日至7月7日每日各蔬菜品类的补货总量与定价策略，使商超在对应时间段内收益最大。。

问题3：考虑到商超销售空间有限，需要进一步制定单品补货方案。要求每日可售单品数量控制在27-33种之间，且每个单品的订购量不低于2.5千克。根据2023年6月24日至6月30日的可售单品情况，确定2023年7月1日各单品的补货数量和定价策略，在尽量满足消费者对各蔬菜品类需求的同时，实现商超收益最大化。

问题4：为了进一步提高蔬菜补货与定价决策的准确性，给出商超应额外采集哪些相关经营数据，并说明这些数据对解决问题有何帮助，给出相应理由。
\section{问题分析}

\subsection{问题一的分析——基于两部分分布模型与多维相关分析的销量规律研究}
针对问题一，需要探究不同蔬菜品类以及单品之间的销售量分布规律与相互关系。
在销售规律方面，由于蔬菜销量中可能出现零销量情况与明显特征，普通简单统计指标难以完成分布规律描述。因此选用两部分分布，通过将各品类日销量划分为"零销量"与"正销量"两部分，通过AIC选取合适的分布规律，通过KS检验判断拟合效果，确定分布规律。
接着，进一步研究不同品类及单品之间是否存在相互关系。
考虑到销售数据具有明显的时间序列特性，直接计算相关系数易受季节变化等共同因素影响，从而产生伪相关，所以需要先对销量进行时间标准化处理，再从不同角度描述相关关系。首先，认定月度季节指数为蔬菜季节性指标；接着在各季节内通过对数变换去除季节影响后计算Sperman相关分析同一季节内销量变化趋势；
接着将日销售量转换为销售份额，减小客流造成的整体性变化，同时利用活跃指数完成蔬菜是否同期销售的判断。最终结合 Bootstrap 置信区间和 FDR 校正对结果进行稳健性检验，得出结论。

\subsection{问题二的分析——基于价格响应与风险优化的补货定价决策}

针对问题二，需要探究在以蔬菜品类为单位进行补货的基础上分析销售总量与成本加成定价之间的关系，并预测2023年7月1日至7月7日各品类每日的补货量以及销售价格使商超收益最大化。
由于蔬菜的销售价格受时令、市场等多种因素影响，若直接根据历史价格与销量进行预测容易出现失误判断。因此，首先需要在控制外部因素的基础上判断价格变化对销售量的实际影响。
得到上述对应影响关系后，进一步考虑到商超在凌晨进行进货时处于不明确当日价格的情况中，故想到可以依据历史批发价的变化规律对未来一周各蔬菜品类的进货成本进行预测，并结合成本加成定价的销售策略得到不同成本条件下的销售价格。
同时结合第一问中的分析，考虑各蔬菜品类与单品之间存在一定的相互关系，对未来一周的需求场景进行讨论，最终计算对应场景的收益，通过构建以期望收益最大化和风险最小化的优化模型，在补货量、销售价格以及库存等条件约束下确定销售策略。

\subsection{问题三的分析——基于需求分解与多目标优化的单品补货决策}

针对问题三，需要在问题二的基础上确定7月1日单品进货策略。
对于单品样本，直接选取最近一周的销量推测容易出现过拟合，完全采用三年平均又会忽略最新状态；因此需要在近期销量、历史可比星期六结构和单品样本量之间建立动态收缩机制，并保留单品份额的不确定性。

接着需要对单品的价格、种类、进货量进行决策。
考虑到题目要求可售单品总数控制在 $27$--$33$ 个，分别固定可售单品数量
$K=27,28,\ldots,33$ ，本文采取两阶段词典序结构作为优化目标：第一阶段优先降低六品类等权的标准化缺货率，第二阶段在服务水平基本不下降的条件下优化期望利润与尾部利润。并比较各方案在全部需求场景下的综合表现。
最终综合考虑市场需求满足程度、收益水平及风险指标，确定最优的单品数量及对应的
补货和定价方案。最终对模型进行检验与灵敏度分析。

\subsection{问题四的分析}

针对问题四，建议商超重点采集两类数据以提升补货与定价决策的准确性。一方面，多品类销量同涨同跌可能受客流总量变化驱动而非品类间的真实替代关系，采集每日门店客流量可在问题一的相关关系分析中排除共同因素干扰，提高品类相关关系判断的准确性。另一方面，蔬菜销量受天气和节假日影响显著，补充历史同期气象数据与节假日标记，可改进问题二中需求预测的季节性调整和问题三中单品份额的动态估计。

\section{模型假设}

\begin{enumerate}[leftmargin=2em,itemsep=0pt,parsep=0pt,topsep=2pt]
    \item 商超的补货决策不影响批发价格，即批发价格由市场供需决定，与单个商超的订货量无关。
    \item 蔬菜的损耗仅与时间相关，与补货量、陈列方式无关，即单品的损耗率在给定时间窗口内为常数。
    \item 历史销售数据中的零销量记录包含真实零需求与缺货导致的未满足需求，且二者的比例在不同品类间相对稳定。
    \item 未来一周（2023 年 7 月 1—7 日）的外部经营环境（天气、节假日、市场竞争等）与历史同期不存在系统性差异。
\end{enumerate}

\needspace{8\baselineskip}
\section{符号说明}
\begin{longtable}{p{2.5em}>{\centering\arraybackslash}p{\dimexpr\textwidth-5em\relax}p{2em}}

\toprule[1.5pt]
符号 & 含义 & 单位 \\
\midrule[1pt]
\endfirsthead
\multicolumn{3}{c}{\tablename~\thetable{}（续表）} \\
\midrule[1pt]
符号 & 含义 & 单位 \\
\midrule[1pt]
\endhead
\bottomrule[1.5pt]
\endfoot
\bottomrule[1.5pt]
\endlastfoot
$c$ & 蔬菜品类编号，$c=1,\ldots,6$ & --- \\
$i$ & 单品编号 & --- \\
$t$ & 日期索引 & 天 \\
$Y_{ct}$ & 品类 $c$ 在第 $t$ 天的总销售量 & kg \\
$Y_{it}$ & 单品 $i$ 在第 $t$ 天的销售量 & kg \\
$\rho_{cc'}$ & 品类 $c$ 与 $c'$ 的 Pearson 相关系数 & --- \\
$J_{cc'}$ & 品类 $c$ 与 $c'$ 的 Jaccard 重合度 & --- \\
$\varepsilon_{ct}$ & 品类 $c$ 在第 $t$ 天的价格弹性 & --- \\
$P_{ct}$ & 品类 $c$ 在第 $t$ 天的零售定价 & 元/kg \\
$W_{ct}$ & 品类 $c$ 在第 $t$ 天的批发价格 & 元/kg \\
$Q_{ct}$ & 品类 $c$ 在第 $t$ 天的补货量 & kg \\
$\hat{D}_{ct}$ & 品类 $c$ 在第 $t$ 天的需求预测值 & kg \\
$\mu_{ct}$ & 品类 $c$ 在第 $t$ 天的成本加成率 & 元/kg \\
$\pi_{ct}$ & 品类 $c$ 在第 $t$ 天的经营利润 & 元 \\
$\lambda_{it}$ & 单品 $i$ 在第 $t$ 天的需求份额 & --- \\
$K$ & 每日可售单品数量 & 个 \\
$\gamma$ & CVaR 风险权重参数 & --- \\
$\text{CVaR}_{\gamma}$ & 均值—下尾 CVaR 风险度量 & 元 \\
$\alpha$ & 缺货惩罚权重 & --- \\
$\beta$ & 利润优化权重 & --- \\
$S_c$ & 品类 $c$ 的标准化缺货率 & --- \\
$h$ & 历史份额权重（长程先验） & --- \\
\end{longtable}

\section{模型建立与求解}
    \subsection{\textbf{数据预处理}}


    \begin{enumerate}[leftmargin=2em]
        \item \textbf{数据清洗}  
        
        \hspace*{2em} 本文首先对所有数据进行清洗处理，发现存在部分缺失值。对缺失值进行合理推测后，认定由于疫情与过年期间的特殊因素导致缺失值出现，并对相关数据进行填0操作，确保数据的完整性与可分析性。

        \item \textbf{数据整合}
        
        \hspace*{2em} 将销售流水数据和批发价格数据按照日期、蔬菜品类和单品进行整合，对单位进行统一，形成一致的数据类型，便于后续分析。
      
    \end{enumerate}
\subsection{问题一的模型建立与求解}
\subsubsection{问题一的数据预处理}
原始数据中，部分商品销售量由于时令影响出现零销售。若直接进行数据分析，易导致相关关系失真。
因此，考虑到数据可靠性，本文采用季节适配条件进行筛选:将单品销售时间跨度不少于730天、正销量天数不少于90天、累计销量不少于100kg的蔬菜纳入考虑范围，筛选后共保留64个单品，覆盖全部蔬菜品类，且均具有真实有效性。
\subsubsection{问题一的建模求解}

\begin{enumerate}[leftmargin=2em,itemsep=0pt,parsep=0pt,topsep=2pt]
    \item \textbf{销售量分布规律的求解}
    
    \hspace*{2em}设第$c$ 类蔬菜在第 $t$ 天的总销售量为 $Y_{ct}$，则有
\begin{equation}
    Y_{ct}=\sum_{i\in I_c}Y_{it}
\end{equation}

其中，$I_c$ 表示属于第$c$类蔬菜的全部单品集合，$Y_{it}$ 表示单品 $i$ 在第 $t$ 天的销售量。
\hspace*{2em}对数据进行处理后，发现各品类蔬菜中存在少量零销量，正销量部分则存在明显右偏与长尾特征。
因此，均值等概率分布指标难以准确描述全部数据，故建立\textbf{两部模型}对日品类销量进行描述，通过将数据分为零销量概率与正销量条件分布来刻画分布规律。
\begin{enumerate}
    \item \textbf{第一部：零销量概率估计} 

\hspace*{2em}记商品 $i$（品类或单品）在 $n_i$ 个观测日中的销量序列为 $\{Y_{i,t}\}_{t=1}^{n_i}$，零销量概率为：

\begin{equation}
 \hat{p}_{0}^{(i)} = \frac{1}{n_i} \sum_{t=1}^{n_i} \mathbf{1}(Y_{it} = 0) \tag{1}
\end{equation}
    \item \textbf{第二部：正销量条件分布拟合}
    
\hspace*{2em}提取正销量子集 $\mathcal{Y}_i^{+} = \{Y_{i,t} : Y_{i,t} > 0\}$。
为避免样本量过小造成参数估计不稳定，对正销量样本数量进行筛选。若 $|\mathcal{Y}_i^{+}| < 30$，认定该品类（或单品）有效销售数据不充足，不进行参数拟合。
对筛选过后的数据，本文选取四种分布作为候选：\textbf{正态分布、}\textbf{对数正态（Lognormal）分布、}\textbf{Gamma 分布、}\textbf{Weibull 分布}，基本契合数据正值右偏特点。

\end{enumerate}
\hspace*{2em}经过对数据的初步整理，发现需要在多个正值右偏模型中选取最优解。因此，本文选用极大似然法估计分布参数，并经由AIC对候选分布做出比较，采取AIC较小者作为优选分布；同时通过KS检验验证所选分布的合理性。
AIC 能够兼顾模型拟合效果与复杂度，而 KS 检验能够从经验分布与理论分布的差异角度进一步检验拟合效果\textsuperscript{\cite{akaike1974}}。 

\hspace*{2em}对于KS检验，若 $p_{\text{KS}} \geq 0.05$，表示所选分布合理；否则表示虽然该分布在候选者中相对最优，但绝对拟合仍不充分。

\begin{table}[H]
    \centering

    \begin{tabularx}{\textwidth}{>{\centering\arraybackslash}X
                                >{\centering\arraybackslash}X
                                >{\centering\arraybackslash}X
                                >{\centering\arraybackslash}X
                                >{\centering\arraybackslash}X
                                >{\centering\arraybackslash}X}
        \toprule
        品类 & 均值/kg & 零销率 & 最优分布 & AIC & KS $p$值 \\
        \midrule
        水生根茎类 & 37.1  & 0.91\% & Gamma     & 9896.7  & 0.490 \\
        花叶类     & 181.4 & 0.91\% & Gamma     & 12442.6 & 0.271 \\
        花菜类     & 38.2  & 1.00\% & Gamma     & 9547.1  & 0.334 \\
        茄类       & 20.5  & 4.11\% & Gamma     & 8058.8  & 0.549 \\
        辣椒类     & 83.7  & 0.91\% & Lognormal & 11036.2 & 0.508 \\
        食用菌     & 69.5  & 0.91\% & Lognormal & 10840.1 & 0.524 \\
        \bottomrule
    \end{tabularx}

    \caption{六类蔬菜日销量分布拟合结果}
    \label{tab:category_distribution}
\end{table}

\hspace*{2em}数据结果如表~\ref{tab:category_distribution}。由结果得出，\textbf{品类层面水生根茎类、花叶类、花菜类和茄类的正销量更符合Gamma分布，辣椒类和食用菌更符合对数正态分布。}
且对应KS检验均大于0.05，表明拟合结果具有较好的合理性。

\begin{figure}[H]
    \centering
    \includegraphics[width=\textwidth]{figures/Q1箱线图——蔬菜品类正相关.png}
    \caption{六类蔬菜日销量分布拟合结果}
    \label{fig:distribution}
\end{figure}

\hspace*{2em}对照图~\ref{fig:distribution} 拟合结果，本文得出结论：

\hspace*{2em}\textbf{从品类层面看，}各品类总体销售较为连续，但在销售规模和波动程度上存在明显差异。其中花叶类销售水平最高且相对稳定，水生根茎类和食用菌的销量波动较强，茄类整体销售水平较低。总体而言，品类层表现为“连续销售、右偏分布、规模分层”的特征。

\hspace*{2em}\textbf{从单品层面看，}不同商品的销售规律差异更加显著。64 个单品的零销率在 1.7%～90.9% 之间变化，表明单品销售覆盖了从常年高频销售到高度季节性销售的多种模式。部分单品全年销售较为连续，而部分商品仅在特定季节集中上市，导致大量零销量日期。同时，不同单品的正销量分布形态也存在明显差异，说明单品层具有较强的个体性。

\hspace*{2em}因此，\textbf{蔬菜品类间聚合后表现较稳定，而单品层则呈现较强的季节性和个体差异。}




\item \textbf{各品类相互关系的求解}


\hspace*{2em}蔬菜具有明显的季节性，不同单品的活跃程度存在较大差异，简单的相关分析不能较好刻画相关关系。
为确保不同年份之间有可比性，将每年7月至次年6月认定为一完整销售年度，并按照春季（3—5 月）、夏季（6—8 月）、秋季（9—11 月）、冬季（12—2 月）四个季节划分。数据覆盖 2020-07 至 2023-06 共三个完整销售年度。

\begin{enumerate}

   \item \textbf{进行K—means聚类}

\hspace*{2em}由于单品种类繁多，互相之间难以直接看出相关关系，本文考虑采用K-means聚类完成相互关系的分析。

\hspace*{2em}\textbf{（1）特征向量构造：} 每个单品以 12 个月销量占比与 12 个月活跃率构成 24 维季节特征向量：
\begin{equation}
        \mathbf{X}_i = (P_{i,1},\ldots,P_{i,12},\ A_{i,1},\ldots,A_{i,12})^\top
 \end{equation}
销量总规模不进入特征，以避免高销量单品仅因规模相近而聚为一类。

\hspace*{2em}\textbf{（2）进行Z-score 标准化：} 占比与活跃率量纲不同，标准化以消除量纲影响、使各维度等权：
\begin{equation}
    x_{i,j}^{*} = \frac{x_{i,j}-\bar{x}_j}{\sigma_j},\qquad j=1,\ldots,24
\end{equation}

\hspace*{2em}\textbf{（3）确定聚类目标：} K-means 通过最小化簇内平方和（WCSS）确定最优划分：
\begin{equation}
    \min_{\{\mathcal{C}_c\}_{c=1}^{k}}\ \sum_{c=1}^{k}\sum_{i\in\mathcal{C}_c}\lVert \mathbf{x}_i^{*} - \boldsymbol{\mu}_c \rVert^2
\end{equation}
其中 $\mathcal{C}_c$ 为第 $c$ 个簇，$\boldsymbol{\mu}_c$ 为簇中心。采用 K-means++ 初始化，对每个 $k$ 重复 100 次随机初始化并取 WCSS 最小者，固定随机种子 $20230907$。

\hspace*{2em}\textbf{（4）选取簇数 $k$ ：} K-means 需预先指定 $k$。对 $k=2,\ldots,8$ 分别聚类，综合轮廓系数、Calinski-Harabasz（CH）指数、Davies-Bouldin（DB）指数、年度块 Bootstrap 稳定性（调整兰德指数 ARI）及最小簇规模约束，选择综合最优的 $k$。计算结果如表 \ref{tab:k} 与图\ref{fig:distribution}所示。
\begin{figure}[H]
    \centering
    \includegraphics[width=\textwidth]{figures/Q1聚类数比较.png}
    \caption{K-means 聚类 $k$ 值选取综合指标}
    \label{fig:K-means}
\end{figure}


\begin{table}[H]
\centering
\begin{tabularx}{\textwidth}{>{\centering\arraybackslash}X>{\centering\arraybackslash}X>{\centering\arraybackslash}X>{\centering\arraybackslash}X>{\centering\arraybackslash}X>{\centering\arraybackslash}X>{\centering\arraybackslash}X}
\toprule
$k$ & 轮廓系数 & CH & DB & 最小簇 & ARI均值 & ARI$\sigma$ \\
\midrule
2 & 0.205 & 19.04 & 1.627 & 27 & 0.593 & 0.244 \\
3 & 0.221 & 18.21 & 1.626 & 15 & 0.540 & 0.346 \\
4 & 0.214 & 16.31 & 1.564 & 10 & 0.585 & 0.279 \\
\textbf{5} & \textbf{0.233} & \textbf{15.79} & \textbf{1.475} & \textbf{7} & \textbf{0.566} & \textbf{0.249} \\
6 & 0.239 & 14.77 & 1.343 & 5 & 0.541 & 0.227 \\
7 & 0.213 & 14.13 & 1.401 & 5 & 0.493 & 0.239 \\
8 & 0.216 & 13.43 & 1.220 & 1 & 0.474 & 0.225 \\
\bottomrule
\end{tabularx}
\caption{K-means 聚类 $k$ 值选取综合指标}
\label{tab:k}
\end{table}

\hspace*{2em}由表 \ref{tab:k}：$k=5$ 时轮廓系数（0.233）与 $k=6$（0.239）接近但 Bootstrap 稳定性更高（ARI 0.566 $>$ 0.541），且最小簇规模（7）满足不小于 5 的约束，季节解释清晰，故确定 $k=5$。聚类结果如表 \ref{tab:cluster} 所示。

\begin{table}[H]
\centering
\begin{tabularx}{\textwidth}{>{\centering\arraybackslash}X>{\centering\arraybackslash}X>{\centering\arraybackslash}X>{\centering\arraybackslash}X>{\centering\arraybackslash}X}
\toprule
簇名 & 规模 & 峰值月 & 主要销售季节 & 代表单品 \\
\midrule
初夏型 & 11 & 6 月 & 夏季 & 木耳菜、苋菜 \\
冬季型 & 7 & 12 月 & 冬季 & 牛肝菌、萝卜 \\
常年型 & 21 & 10 月 & 四季均匀 & 菠菜、大白菜 \\
春季型 & 8 & 3 月 & 春季 & 蒜苗、红苋菜 \\
夏秋型 & 17 & 9 月 & 夏秋 & 白菜、金针菇 \\
\bottomrule
\end{tabularx}
\caption{五簇销售模式画像}
\label{tab:cluster}
\end{table}



\hspace*{2em}聚类结果表明：64 个单品可归纳为初夏型、冬季型、常年型、春季型、夏秋型五种季节销售模式，其中常年型占 33\%，说明相当一部分蔬菜已实现全年供应。聚类将"品类—单品"两级结构拓展为"品类—单品—季节簇"三级结构，为分层相关分析提供分组依据。

 

    
    \item \textbf{构造单品季节特征}
  
    \textbf{（1）月均日销量 $\bar{Q}_{i,m}$} 用于消除不同月份天数差异对销量比较的影响：
\begin{equation}
    \bar{Q}_{i,m} = \frac{1}{N_y}\sum_{y=1}^{N_y}\frac{Q_{i,y,m}}{d_m}
\end{equation}

\textbf{（2）月度销量占比 $P_{i,m}$} 用于刻画该月销量占全年的比重，反映需求的月份集中程度：
\begin{equation}
    P_{i,m} = \frac{1}{N_y}\sum_{y=1}^{N_y}\frac{Q_{i,y,m}}{\displaystyle\sum_{m'=1}^{12}Q_{i,y,m'}}
\end{equation}

\textbf{（3）月度销售活跃率 $A_{i,m}$} 用于刻画该月有实际销售的日期占比，可区分"旺季销量高"与"仅在旺季上市"两种现象：
\begin{equation}
    A_{i,m} = \frac{1}{N_y}\sum_{y=1}^{N_y}\frac{\displaystyle\sum_{t\in\mathcal{T}_{i,y,m}}\mathbf{1}(Y_{i,t}>0)}{d_m}
\end{equation}

\textbf{（4）季节指数 $S_{i,m}$} 用于刻画该月日均销量相对全年平均水平的强弱，$S_{i,m}>1$ 表示高于全年平均，$S_{i,m}<1$ 表示低于全年平均：
\begin{equation}
    S_{i,m} = \frac{\bar{Q}_{i,m}}{\displaystyle\frac{1}{12}\sum_{m'=1}^{12}\bar{Q}_{i,m'}}
\end{equation}

\hspace*{2em}上述四个指标分别从规模、结构、频次、相对强弱四个角度刻画季节特征，为后续数据处理提供统一的量化基础。得到结果如图~\ref{fig:season}所示。
    
\begin{figure}[H]
    \centering
    \includegraphics[width=\textwidth]{figures/fig_q1_cluster_profile_heatmaps.png}
    \caption{五个季节型簇的季节指数与销售活跃率双画像}
    \label{fig:season}
\end{figure}



    \item \textbf{多指标分析相关关系}
    
    \hspace*{2em}完成上述分析后，需要进一步对不同品类（或单品）之间研究相关关系。为此，本文构建了四个指标进行衡量。

    \textbf{（1）季节指数 Pearson 相关 $\rho_{ij}^{\mathrm{SI}}$} 
    
    \hspace*{2em}以季节指数衡量两实体旺淡季曲线的同步程度：
\begin{equation}
    \rho_{ij}^{\mathrm{SI}} = \frac{\displaystyle\sum_{m=1}^{12}\big(S_{i,m}-\bar{S}_i\big)\big(S_{j,m}-\bar{S}_j\big)}
    {\sqrt{\displaystyle\sum_{m=1}^{12}\big(S_{i,m}-\bar{S}_i\big)^2}\ \sqrt{\displaystyle\sum_{m=1}^{12}\big(S_{j,m}-\bar{S}_j\big)^2}}
\end{equation}
\hspace*{2em}正值表示旺季月份接近（同步），负值表示旺淡季错位（互补）。


\textbf{（2）分季节销量 Spearman 相关 $\rho_{ij}^{\mathcal{S}}$} 

\hspace*{2em}用于衡量同一季节内两实体日销量的同向波动。为压缩右偏并消除不同月份基准水平差异，先对日销量做对数变换并减去所属年月均值：
\begin{equation}
    \tilde{Y}_{i,t} = \log(1+Y_{i,t}) - \frac{1}{|\mathcal{T}_{i,y,m}|}\sum_{t'\in\mathcal{T}_{i,y,m}}\log(1+Y_{i,t'})
\end{equation}
\hspace*{2em}再在各季节 $\mathcal{S}$ 内计算去趋势序列的 Spearman 秩相关系数：
\begin{equation}
    \rho_{ij}^{\mathcal{S}} = 1 - \frac{6\sum_{t\in\mathcal{S}}d_{ij,t}^{2}}{n_{\mathcal{S}}(n_{\mathcal{S}}^{2}-1)}
\end{equation}
其中 $d_{ij,t}$ 为 $\tilde{Y}_{i,t}$ 与 $\tilde{Y}_{j,t}$ 在季节 $\mathcal{S}$ 内的秩差，$n_{\mathcal{S}}$ 为该季节天数。

\textbf{（3）分季节销量份额 Spearman 相关} 

\hspace*{2em}用于剥离共同客流造成的同时上涨。将每日销量转换为当日份额：
\begin{equation}
    R_{i,t}^{\mathcal{S}} = \frac{Y_{i,t}}{\sum_{j\in\text{level}}Y_{j,t}},\qquad t\in\mathcal{S}
\end{equation}
\hspace*{2em}再按季节计算份额序列的 Spearman 相关，负值提示销售份额竞争或错位。

\textbf{（4）活跃日 Jaccard 重合度 $J_{ij}$} 

\hspace*{2em}用于衡量两实体经常同期有销售的活跃日集合的重合程度：
\begin{equation}
    J_{ij} = \frac{|A_i \cap A_j|}{|A_i \cup A_j|}
\end{equation}
其中 $A_i=\{t:Y_{i,t}>0\}$ 为正销售日集合，$J_{ij}$ 越高表示同期销售程度越高。

\hspace*{2em}由于日销量序列存在时间依赖，采用 7 日移动块 Bootstrap（200 次重采样）估计 95\% 置信区间进行可信度估量，并采用 Benjamini-Hochberg（BH）校正控制错误发现率（FDR），避免多重比较导致的假阳性。

\end{enumerate}
\end{enumerate}

\hspace*{2em}最终得到结果如图~\ref{fig:Person}（见附录）与图~\ref{fig:Sperman}所示。

\begin{figure}[H]
    \centering
    \includegraphics[width=\textwidth]{figures/fig_q1_hierarchical_relationships.png}
    \caption{ 品类与季节簇的总量共变及份额竞争矩阵}
    \label{fig:Sperman}
\end{figure}
 
故得到\textbf{结论：}
\begin{enumerate}
    \item \textbf{品类间关系（15 对）}
    
 \hspace*{2em}四种指标的计算结果汇总于表 \ref{tab:cat_rel}。

\begin{table}[H]
\centering
\begin{tabularx}{\textwidth}{>{\hsize=0.92\hsize\centering\arraybackslash}X>{\hsize=1.0\hsize\centering\arraybackslash}X>{\hsize=1.4\hsize\centering\arraybackslash}X>{\hsize=0.6\hsize\centering\arraybackslash}X>{\hsize=1.08\hsize\centering\arraybackslash}X}
\toprule
关系指标 & 取值范围 & 最强对 & 关键值 & 最弱对 \\
\midrule
全年销量相关 & $[0.415,0.692]$ & \textbf{花叶类—辣椒类} & 0.692 & 花菜类—茄类 \\
冬季销量相关 & $[0.564,0.791]$ & \textbf{整体最强} & —— & —— \\
夏季销量相关 & $[0.222,0.617]$ & —— & —— & 分化最大 \\
季节指数相关 & $[-0.677,0.841]$ & \textbf{水生根茎类—食用菌} & 0.841 & 茄类—食用菌 \\
全年份额相关 & 最低 $-0.535$ & 花叶类—食用菌 & —— & —— \\
活跃日 & $[0.967,1.000]$ & 几乎每日共售 & —— & —— \\
\bottomrule
\end{tabularx}
\caption{品类间关系指标汇总}
\label{tab:cat_rel}
\end{table}



\hspace*{2em}
由表 \ref{tab:cat_rel} 并结合图 \ref{fig:Person}、图 \ref{fig:Sperman} 可见：
品类全年销量相关均为正（$[0.415,0.692]$），最强为花叶类—辣椒类，冬季最强、夏季分化最大；但全年份额相关可降至 $-0.535$，说明总量同涨与结构此消彼长并存，总量相关与份额相关不可互相替代。

\item \textbf{簇间关系（10 对）}


 \hspace*{2em}五簇聚合销量后的关系结果汇总于表 \ref{tab:clu_rel}，对应矩阵见图 \ref{fig:Sperman}。

\begin{table}[H]
\centering
\begin{tabularx}{\textwidth}{>{\hsize=0.92\hsize\centering\arraybackslash}X>{\hsize=1.0\hsize\centering\arraybackslash}X>{\hsize=1.4\hsize\centering\arraybackslash}X>{\hsize=0.6\hsize\centering\arraybackslash}X>{\hsize=1.08\hsize\centering\arraybackslash}X}
\toprule
关系指标 & 数值 & 关键对 & 关键值 & 说明 \\
\midrule
全年销量相关 & \textbf{$[-0.018,0.393]$} & \textbf{常年型—冬季型} & 0.393 & 明显弱于品类层 \\
季节指数相关 & \textbf{$-0.734$} & 春季型—夏秋型 & —— & 峰值月份错位 \\
全年份额相关 & \textbf{$-0.787$} & 初夏型—冬季型 & —— & 结构反向最强 \\
\bottomrule
\end{tabularx}
\caption{簇间关系指标汇总}
\label{tab:clu_rel}
\end{table}

\hspace*{2em}由表 \ref{tab:clu_rel} 可见：
簇间销量相关（$[-0.018,0.393]$）显著弱于品类层，说明聚类有效分离了不同销售模式；常年型—冬季型相关最高，春季型—夏秋型呈峰值错位，初夏型—冬季型份额相关为 $-0.787$，但仅解释为层级内组成变化。



\item \textbf{单品间关系（2016 对）}

 \hspace*{2em}关系强度分布如表 \ref{tab:sku_dist} 所示，代表性单品对的销量与季节关系如图 \ref{fig:sku_pairs} 所示，两对代表性关系的完整指标列于表 \ref{tab:sku_rep}。

\begin{table}[H]
\centering
\begin{tabularx}{\textwidth}{>{\centering\arraybackslash}X>{\centering\arraybackslash}X>{\centering\arraybackslash}X}
\toprule
强度级别 & 对数 & 占比 \\
\midrule
\textbf{较明确关系} & \textbf{18} & \textbf{0.89\%} \\
弱至中等关系 & 221 & 11.0\% \\
仅保留 & 1777 & 88.1\% \\
\bottomrule
\end{tabularx}
\caption{单品关系强度分布}
\label{tab:sku_dist}
\end{table}


\begin{figure}[H]
\centering
\includegraphics[width=\textwidth]{figures/fig_q1_representative_pairs.png}
\caption{代表性单品对的销量与季节关系}
\label{fig:sku_pairs}
\end{figure}


\begin{table}[H]
\centering
\begin{tabularx}{\textwidth}{>{\hsize=1.8\hsize\centering\arraybackslash}X>{\hsize=0.6\hsize\centering\arraybackslash}X>{\hsize=1.2\hsize\centering\arraybackslash}X>{\hsize=0.7\hsize\centering\arraybackslash}X>{\hsize=0.7\hsize\centering\arraybackslash}X}
\toprule
单品对 & 全年 $\rho$ & 95\% CI & 季节指数 $r$ & Jaccard \\
\midrule
西兰花—净藕(1) & \textbf{0.442} & $0.381$～$0.515$ & 0.549 & \textbf{0.974} \\
西峡花菇(1)—西峡香菇(1) & \textbf{$-0.326$} & $-0.435$～$-0.227$ & $-0.619$ & 0.061 \\
\bottomrule
\end{tabularx}
\caption{代表性单品对关系指标}
\label{tab:sku_rep}
\end{table}
\hspace*{2em}由表 \ref{tab:sku_dist} 可见：
2016 对中仅 18 对（0.89\%）为较明确关系（17 正 1 负），221 对为弱至中等，其余 1777 对（88.1\%）无显著关联。由表 \ref{tab:sku_rep} 并结合图 \ref{fig:sku_pairs}：最高相关为西兰花与净藕(1)（0.442），唯一明确负相关为西峡花菇(1)与西峡香菇(1)（$-0.326$），主要体现上市期错位。单品层关系显著弱于品类层，统计上稳定的共变关系稀少。




\end{enumerate}





\subsection{问题二的模型建立与求解} 

\subsubsection{问题二的建模求解过程}

\begin{enumerate}

    \item \textbf{建立决策变量与数据口径}

\hspace*{2em}设品类集合 $\mathcal{C}=\{$水生根茎类、花叶类、花菜类、茄类、辣椒类、食用菌$\}$，决策日集合 $\mathcal{T}=\{2023\text{-}07\text{-}01,\ldots,07\text{-}07\}$。决策变量为各日各品类的损耗前补货量 $q_{ct}$ 与成本加成率 $m_{ct}$，建议售价由预测成本与加成率联合确定：
\begin{equation}
    p_{ct} = \hat{W}_{ct}\big(1+m_{ct}\big),\qquad c\in\mathcal{C},\ t\in\mathcal{T}
\end{equation}

历史销量 $Q_{ct}$ 取品类日毛销售量，售价 $P_{ct}$ 取销量加权均价，批发成本 $W_{ct}$ 取当日已售单品批发价的销量加权均值，近期损耗率 $\lambda_c$ 取最近 30 天的销量加权值。

   \item \textbf{建立价格--销量关系模型}

\hspace*{2em}蔬菜销量不仅受到销售价格的影响，还与季节、星期、节假日以及历史销售情况等因素有关。
若直接对价格与销量进行回归，其他因素可能对价格效应的估计产生干扰。因此，本文采用部分线性模型，
将销售价格作为主要解释变量，其余影响因素作为控制变量，建立价格与销量之间的关系：

\begin{equation}
    Y_{ct}=\theta_c X_{ct}+g_c(Z_{ct})+u_{ct},
    \qquad
    X_{ct}=m_c(Z_{ct})+v_{ct}.
\end{equation}

其中，$Y_{ct}=\ln(Q_{ct}+1)$ 为第 $c$ 类蔬菜第 $t$ 日的对数销量，
$X_{ct}=\ln P_{ct}$ 为对数销售价格，$\theta_c$ 为价格响应系数，
$Z_{ct}$ 为控制变量，包括星期、月份、时间趋势、当期成本以及历史销量和价格等信息。
$g_c(\cdot)$ 和 $m_c(\cdot)$ 用于描述控制变量与销量、价格之间的关系，
$u_{ct}$ 和 $v_{ct}$ 为随机误差项。


\hspace*{2em}接着分三步。第一步，用梯度提升树（GBDT）分别拟合 ${E}[Y|Z]$ 和 ${E}[X|Z]$，将估计值从原始数据中去除，得到剔除了混杂因素后的残差：
\begin{equation}
    \tilde{Y}_{ct} = Y_{ct} - \hat{g}_c(Z_{ct}), \qquad \tilde{X}_{ct} = X_{ct} - \hat{m}_c(Z_{ct}).
\end{equation}

\hspace*{2em}第二步，对 $\tilde{Y}_{ct}$ 关于 $\tilde{X}_{ct}$ 做普通最小二乘回归，斜率即为价格弹性 $\hat{\theta}_c$。

\hspace*{2em}第三步，用五折时间序列交叉验证确保训练集严格早于测试集，避免信息泄露。

用 7 日移动窗口的块自举法（Bootstrap）计算弹性系数的标准误和 95\% 置信区间。部分品类的价格变化主要来自商超主动调价，原始估计可能不可靠，因此引入公共收缩：
若某品类的弹性估计为正，表明价格与销量呈正向关系，与需求理论中价格与需求量通常呈反向变动的基本规律不一致\textsuperscript{\cite{marshall1920}}；同时，已有实证研究表明，不同品类及估计方法可能导致价格弹性估计存在较大差异\textsuperscript{\cite{tellis1988}}。因此，对于此类异常估计值，本文采用全品类合并后的公共弹性进行收缩替代。最终采用的弹性系数为
\begin{equation}
    \tilde{\theta}_c = \mathrm{Proj}_{[-3,-0.05]}\Big\{\omega_c\hat{\theta}_c + (1-\omega_c)\hat{\theta}_0\Big\}
\end{equation}
其中 $\hat{\theta}_0$ 为全品类合并弹性，$\omega_c$ 为经验贝叶斯权重，投影操作 $\mathrm{Proj}_{[-3,-0.05]}$ 将结果限制在合理的区间内（弹性绝对值不超过 3、不低于 0.05）。

\hspace*{2em}一旦知道了每个品类的价格弹性，就可以对任意候选售价 $p$，精确计算出它能把需求提升或抑制多少。




\hspace*{2em}六品类价格弹性的估计、95\% 置信区间与优化采用值如表 \ref{tab:q2_resp} 所示，对应图形见图 \ref{fig:q2_response}。

\begin{table}[H]
\centering
\begin{tabularx}{\textwidth}{>{\hsize=1.0\hsize\centering\arraybackslash}X>{\hsize=0.8\hsize\centering\arraybackslash}X>{\hsize=1.6\hsize\centering\arraybackslash}X>{\hsize=0.8\hsize\centering\arraybackslash}X>{\hsize=0.8\hsize\centering\arraybackslash}X}
\toprule
品类 & 原始 $\theta$ & 95\% 块 Bootstrap 区间 & 优化用 $\theta$ & 合并回退 \\
\midrule
水生根茎类 & $-1.858$ & $[-2.344,-1.335]$ & \textbf{$-1.688$} & 否 \\
花叶类 & $-0.161$ & $[-0.878,0.609]$ & $-0.181$ & 否 \\
花菜类 & $-1.009$ & $[-1.595,-0.202]$ & \textbf{$-0.861$} & 否 \\
茄类 & $0.270$ & $[-0.698,1.073]$ & $-0.246$ & \textbf{是} \\
辣椒类 & $0.061$ & $[-0.140,0.244]$ & $-0.246$ & \textbf{是} \\
食用菌 & $-0.246$ & $[-0.935,0.369]$ & $-0.246$ & 否 \\
\bottomrule
\end{tabularx}
\caption{品类价格弹性估计}
\label{tab:q2_resp}
\end{table}

\begin{figure}[H]
\centering
\includegraphics[width=0.85\textwidth]{figures/fig_q2_price_response_forest.png}
\caption{价格弹性估计与优化采用值}
\label{fig:q2_response}
\end{figure}

\hspace*{2em}
由表 \ref{tab:q2_resp}：\textbf{水生根茎类和花菜类的价格响应最显著}（优化值约 $-1.69$ 和 $-0.86$），模型对这两类的定价判断较为可靠；\textbf{茄类和辣椒类无法识别出显著的价格弹性}，模型回退至全品类合并响应 $-0.246$ 并收紧加价上界。该表支持"控制可观测状态后的条件价格响应"解释，六个系数不宜全部视为显著因果弹性。

\item \textbf{建立成本预测模型}

\hspace*{2em}补货决策需要在凌晨做出，但当日的实际进货成本在决策时点尚未完全可知。若成本预测偏差较大，基于不准确成本计算的补货量和定价都会失效，因此需要提前预测未来一周各品类每日的进货成本。

\hspace*{2em}成本预测同时运行两个模型并根据历史表现动态分配权重。第一个是梯度提升树（GBDT）模型，输入包括星期、月份、历史成本滞后项、品类销量信息等，用于捕获非线性规律。第二个是季节朴素模型——直接用上年同期的成本作为预测值，作为一个简单但稳健的基准。在历史数据上做五折时间序列验证，记录每个模型在每一折的预测误差。若 GBDT 在某品类上持续优于朴素模型，就赋予它更大的权重；若数据不足导致 GBDT 过拟合，权重自动回退到朴素模型。设 GBDT 模型和季节朴素模型在验证集上的预测均方误差分别为 $\mathrm{MSE}_c^{\mathrm{GBDT}}$ 和 $\mathrm{MSE}_c^{\mathrm{naive}}$，则 GBDT 的权重为
\begin{equation}
    w_c = \max\left\{0,\ 1 - \frac{\mathrm{MSE}_c^{\mathrm{GBDT}}}{\mathrm{MSE}_c^{\mathrm{naive}}}\right\},
\end{equation}
\hspace*{2em}即 GBDT 误差低于朴素基准时赋予正权重，否则权重为 0、完全回退到朴素模型。最终预测进货成本为两者的加权组合：
\begin{equation}
    \hat{W}_{ct} = w_c \cdot \hat{W}_{ct}^{\mathrm{GBDT}} + (1-w_c) \cdot \hat{W}_{ct}^{\mathrm{naive}}.
\end{equation}

\hspace*{2em}五折递增式时间验证结果如表 \ref{tab:q2_forecast} 所示。


\begin{table}[H]
\centering
\begin{tabularx}{\textwidth}{>{\centering\arraybackslash}X>{\centering\arraybackslash}X>{\centering\arraybackslash}X>{\centering\arraybackslash}X>{\centering\arraybackslash}X>{\centering\arraybackslash}X}
\toprule
品类 & GBDT权重 & 需求WAPE & 朴素WAPE & 成本WAPE & 80\%覆盖 \\
\midrule
水生根茎类 & 0.45 & 44.9\% & 52.6\% & 16.4\% & 73.5\% \\
花叶类 & 0.25 & 27.0\% & 29.5\% & 13.0\% & 74.2\% \\
花菜类 & 0.35 & 38.5\% & 43.5\% & 14.1\% & 76.9\% \\
茄类 & 0.15 & 39.0\% & 40.9\% & 14.4\% & 61.2\% \\
辣椒类 & 0.20 & 31.7\% & 34.6\% & 17.5\% & 80.4\% \\
食用菌 & 0.35 & 34.3\% & 38.2\% & 16.3\% & 80.0\% \\
\bottomrule
\end{tabularx}
\caption{五折滚动时间验证结果}
\label{tab:q2_forecast}
\end{table}

\hspace*{2em}
由表 \ref{tab:q2_forecast}：\textbf{需求预测加权 WAPE 为 35.88\%}，较季节朴素法（39.85\%）下降约 10\%；\textbf{成本预测 WAPE 为 15.28\%}。需求 80\% 经验区间覆盖率 74.36\% 低于名义水平，说明必须保留多情景随机优化——若完全依赖点预测做确定性优化，会系统性低估低需求风险。

\item \textbf{建立联合需求场景模型}

\hspace*{2em}有了价格弹性和成本预测，理论上可以直接预测未来一周每天每个品类的需求量。但预测不可能完全准确。表 \ref{tab:q2_forecast} 显示需求预测的平均误差（WAPE）约为 36\%，区间覆盖率也低于目标值。若基于一个有 36\% 误差的单一数值做补货决策，要么经常缺货、要么经常积压。因此，模型不给出单一预测值，而是生成一组覆盖各种可能性的需求情景。

\hspace*{2em}在做需求预测之前，需要先从历史销量中剔除价格因素的影响，得到在正常价格水平下的潜在需求场景：
\begin{equation}
    A_{ct} = Y_{ct} - \tilde{\theta}_c \ln P_{ct}
\end{equation}
\hspace*{2em}这一操作的目的是把历史销量中由价格波动引起的那部分变动去掉，剩下由季节、星期、市场环境等非价格因素驱动的真实需求水平。对未来这个潜在需求做预测后，对于任意候选售价 $p$，实际需求可以通过弹性系数还原回来：
\begin{equation}
    \hat{D}_{ct}(p) = \Big[\exp\big\{\hat{A}_{ct} + \tilde{\theta}_c \ln p\big\} - 1\Big]_+
\end{equation}
\hspace*{2em}预测时只使用决策日之前就能确定的信息（星期、月份、历史滞后数据等），不使用前一天的实际销量作为输入，避免用未来的信息预测未来。

\hspace*{2em}第一问已证实六品类销量之间存在正相关——市场整体向好时各品类一起好、市场低迷时一起差。若每个品类独立生成需求情景，会人为削弱这种联动效应，导致在糟糕日子里低估多个品类同时滞销的风险。为保留品类间的真实相关性，模型将六个品类的需求预测残差与六个品类的成本预测残差按日期拼接成一个 12 维向量，截断极端值后进行中心化处理：
\begin{equation}
    \mathbf{r}_t = \big(e_{1t}^{D},\ldots,e_{6t}^{D},\ e_{1t}^{W},\ldots,e_{6t}^{W}\big)^{\top}
\end{equation}
\hspace*{2em}然后从历史数据中所有连续的 7 天时间段里随机抽取 600 个（允许重复），每个 7 天时间段产生一组未来一周的联合残差。将抽取的残差加回到对应的需求预测基线 $\hat{D}_{ct}$ 上，得到第 $\omega$ 组情景下的实际需求：
\begin{equation}
    D_{ct}^{(\omega)} = \hat{D}_{ct} \cdot \exp\big(e_{ct}^{(\omega)}\big).
\end{equation}
\hspace*{2em}这样生成的情景同时保留了三个层面的相关性：品类之间的联动效应、成本与需求之间的相关性、以及一周内连续几天之间的时间连续性。




\item \textbf{建立补货—定价联合优化模型}

\hspace*{2em}有了各情景下的需求预测和成本预测，最后一步是确定每天每个品类的补货量 $q_{ct}$ 和定价 $p_{ct}$。但若只追求平均利润最大，模型可能会选出''大部分情景下赚钱但偶尔大幅亏损''的方案。在蔬菜经营中，连续亏损会迅速侵蚀现金流，导致后续连进货都困难。因此优化目标不能只看平均赚多少，还要看最坏情况下能保住多少。

\hspace*{2em}每天每个品类的补货量 $q_{ct}$ 在凌晨确定后，由于损耗率 $\lambda_c$ 的存在，实际可售数量为 $(1-\lambda_c)q_{ct}$。实际销量取需求与可售量中的较小值：
\begin{equation}
    S_{ct}^{(\omega)} = \min\big\{D_{ct}^{(\omega)}(p_{ct}),\ (1-\lambda_c)q_{ct}\big\}
\end{equation}
\hspace*{2em}单组情景下的总利润为各品类各天的销售收入减去进货总成本：
\begin{equation}
    \Pi^{(\omega)} = \sum_t\sum_c\Big[p_{ct} S_{ct}^{(\omega)} - W_{ct}^{(\omega)} q_{ct}\Big]
\end{equation}
\hspace*{2em}设 $g_{ct} = \big[D_{ct}^{(\omega)}(p_{ct}) - S_{ct}^{(\omega)}\big]_+$ 为第 $c$ 品类第 $t$ 日的缺货量，则修正后的服务调整利润为
\begin{equation}
    \Pi_{svc} = \Pi_{op} - \eta \sum_t \sum_c\, g_{ct} \cdot \big[p_{ct} - W_{ct}^{(\omega)}\big]_+,
\end{equation}
\hspace*{2em}其中 $\eta$ 为单位缺货量造成的商誉损失系数。缺货量乘以单品毛利 $[p_{ct} - W_{ct}^{(\omega)}]_+$，是因为只有盈利商品缺货才会造成直接的利润损失和顾客流失。

\hspace*{2em}模型不只看利润的平均水平，还要控制最差情况下的收益。将 600 个情景下的利润从小到大排序，取最差的 10\% 情景的平均值（称为下尾条件风险值，记作 LCVaR）：
\begin{equation}
    \mathrm{LCVaR}_{\beta}(\Pi) = \max_{\eta}\Big\{\eta - \frac{1}{\beta\Omega}\sum_{\omega}\big(\eta - \Pi^{(\omega)}\big)_+\Big\}
\end{equation}
\hspace*{2em}这个值的含义是：在最坏的 10\% 情况中，平均能赚多少钱。值越高说明最坏情况下的收益越有保障。

\hspace*{2em}最终优化目标是在平均利润和下尾利润之间取得平衡：
\begin{equation}
    \max\ J = (1-\gamma)\,{E}\big[\Pi_{svc}\big] + \gamma\,\mathrm{LCVaR}_{0.10}\big(\Pi_{svc}\big)
    - \lambda_{ref}\sum_t\sum_c\big(m_{ct} - m_c^{med}\big)^2
\end{equation}
其中 $\gamma$ 控制决策者对下尾安全的重视程度（主方案取 0.25，即平均利润占 75\%、下尾利润占 25\%），$\lambda_{ref}$ 是一个惩罚项，用来防止定价偏离历史正常范围太多。加成率的上下界取自最近一年的历史分位数，相邻两天的加成率变化幅度也受到限制，避免价格大幅跳变。补货量同样有非负约束和上限约束：
\begin{equation}
    m_c^{L}\leq m_{ct}\leq m_c^{U},\quad \big|m_{ct}-m_{c,t-1}\big|\leq \Delta m_c,\quad 0\leq q_{ct}\leq q_c^{U}
\end{equation}
\hspace*{2em}主方案取 $\gamma=0.25$、$\beta=0.10$、$\lambda_{ref}=80$、商誉系数 $\eta=0.10$；同时求解 $\gamma=0$（纯利润最大化）与 $\gamma=0.50$（偏重安全）作风险对照，并对关键参数作敏感性重优化。
给出的 2023 年 7 月 1--7 日各品类日补货总量与定价策略如图 \ref{fig:q2_strategy}、图 \ref{fig:q2_similarity} 及表 \ref{tab:q2_strategy} 所示，每日汇总见表 \ref{tab:q2_daily}，一周品类汇总见表 \ref{tab:q2_cat}。

\begin{table}[H]
\centering
\begin{tabularx}{\textwidth}{>{\centering\arraybackslash}X>{\centering\arraybackslash}X>{\centering\arraybackslash}X>{\centering\arraybackslash}X}
\toprule
日期 & 补货总量/kg & 场景平均销售/kg & 经营利润期望/元 \\
\midrule
2023-07-01 & 362.26 & 302.82 & 750.54 \\
2023-07-02 & 281.04 & 230.98 & 666.82 \\
2023-07-03 & 241.96 & 199.74 & 585.46 \\
2023-07-04 & 258.61 & 212.91 & 636.63 \\
2023-07-05 & 282.18 & 230.89 & 667.17 \\
2023-07-06 & 310.26 & 254.25 & 767.22 \\
2023-07-07 & 312.02 & 255.18 & 776.45 \\
\bottomrule
\end{tabularx}
\caption{主方案每日汇总}
\label{tab:q2_daily}
\end{table}

\begin{table}[H]
\centering
\begin{tabularx}{\textwidth}{>{\hsize=0.92\hsize\centering\arraybackslash}X>{\hsize=0.78\hsize\centering\arraybackslash}X>{\hsize=1.0\hsize\centering\arraybackslash}X>{\hsize=1.0\hsize\centering\arraybackslash}X>{\hsize=1.0\hsize\centering\arraybackslash}X>{\hsize=1.3\hsize\centering\arraybackslash}X}
\toprule
品类 & 补货/kg & 平均销售/kg & 平均售价/元 & 平均加成率 & 经营利润/元 \\
\midrule
水生根茎 & 65.90 & 50.59 & 20.14 & 65.3\% & 197.47 \\
花叶类 & 900.05 & 739.40 & 6.74 & 98.7\% & 1908.32 \\
花菜类 & 106.64 & 85.16 & 14.94 & 86.3\% & 420.75 \\
茄类 & 123.92 & 101.16 & 8.45 & 87.2\% & 292.81 \\
辣椒类 & 559.76 & 451.85 & 7.79 & 127.4\% & 1528.81 \\
食用菌 & 292.06 & 258.62 & 6.02 & 70.7\% & 502.15 \\
\bottomrule
\end{tabularx}
\caption{一周品类汇总}
\label{tab:q2_cat}
\end{table}
\begin{figure}[H]
\centering
\includegraphics[width=0.85\textwidth]{figures/fig_q2_daily_replenishment_pricing.png}
\caption{未来一周品类级补货与定价策略}
\label{fig:q2_strategy}
\end{figure}

\begin{figure}[H]
\centering
\includegraphics[width=0.85\textwidth]{figures/fig_q2_strategy_heatmaps.png}
\caption{补货与定价策略的日期间相似性}
\label{fig:q2_similarity}
\end{figure}


\hspace*{2em}\textbf{注：}2023 年 7 月 1—7 日各品类逐日补货总量与定价策略详见附录表~\ref{tab:q2_strategy}（见附录）。





\hspace*{2em}
由图 \ref{fig:q2_strategy} 与表 \ref{tab:q2_cat}：\textbf{花叶类和辣椒类是主要利润来源}，与其在第一问中的高销量与高稳定性一致；\textbf{水生根茎类因价格弹性最强（约 $-1.69$），模型采用高加成率（65\%）与较小补货量控制库存风险}。补货量日间波动较大而售价相对平稳（图 \ref{fig:q2_similarity}），说明模型主要通过数量端吸收短期需求波动，价格调整滞后于采购调整。\textbf{主方案一周补货 2048.33 kg，经营利润期望 4850.30 元。}


\item \textbf{收益与风险比较}

\hspace*{2em}主方案相对历史中位基线的收益改善如图 \ref{fig:q2_improvement} 所示，不同风险权重下的期望与下尾表现见表 \ref{tab:q2_risk} 及图 \ref{fig:q2_tradeoff}。

\begin{figure}[H]
\centering
\includegraphics[width=0.85\textwidth]{figures/fig_q2_profit_dumbbell.png}
\caption{主策略相对历史基线的收益改善}
\label{fig:q2_improvement}
\end{figure}

\begin{table}[H]
\centering
\begin{tabularx}{\textwidth}{>{\centering\arraybackslash}X>{\centering\arraybackslash}X>{\centering\arraybackslash}X>{\centering\arraybackslash}X}
\toprule
风险权重 $\gamma$ & 经营利润期望/元 & 正则化利润期望/元 & 正则化下尾利润/元 \\
\midrule
0.00 & 4857.62 & 4184.90 & 205.01 \\
\textbf{0.25} & \textbf{4850.30} & \textbf{4169.48} & \textbf{302.51} \\
0.50 & 4792.95 & 4098.13 & 413.87 \\
\bottomrule
\end{tabularx}
\caption{风险权重敏感性}
\label{tab:q2_risk}
\end{table}

\hspace*{2em}
由图 \ref{fig:q2_improvement}：\textbf{主方案四项指标均优于历史基线}，周经营利润期望由约 2671 元提高至 4850 元，最差 10\% 经营利润由负转正，均值—CVaR 双目标在提升收益的同时改善了尾部安全。由表 \ref{tab:q2_risk}：$\gamma$ 从 0 增至 0.50 时，正则化下尾利润由 205 元提升至 414 元，而期望利润仅由 4185 元降至 4098 元，边际代价递减；\textbf{主方案 $\gamma=0.25$ 位于拐点附近}，以较小期望损失换取明显尾部改善。

\begin{figure}[H]
\centering
\includegraphics[width=0.85\textwidth]{figures/fig_q2_risk_return_tradeoff.png}
\caption{不同风险权重下的期望—下尾权衡}
\label{fig:q2_tradeoff}
\end{figure}




\item \textbf{关键参数敏感性分析验证}


\begin{figure}[H]
\centering
\includegraphics[width=0.85\textwidth]{figures/fig_q2_management_sensitivity.png}
\caption{关键管理参数敏感性分析}
\label{fig:q2_sensitivity}
\end{figure}

\hspace*{2em}
由图 \ref{fig:q2_sensitivity}：\textbf{商誉成本系数越大，补货量越高、缺货概率越低}（系数 0→0.20，缺货概率 66.3\%→57.9\%），说明服务损失惩罚能有效推动提高供给保障；\textbf{参考加价正则权重越大，平均加成率越低、极端定价决策越少}（权重 0→500，加成率 0.90→0.72），说明该参数能抑制极端定价。\textbf{主方案取商誉系数 0.10、正则权重 80}，在服务保障与价格稳定之间取得平衡。

\hspace*{2em}
根据上述内容得出\textbf{结论：}基于部分线性双重机器学习估计六品类条件价格响应，以品类间 Pearson 相关性为依据拼接 12 维联合残差向量，生成 600 个保留跨品类相关性的需求场景，以均值—下尾 CVaR 为优化目标联合求解 42 个补货量与 42 个加成率。\textbf{主方案一周补货 2048.33 kg，经营利润期望 4850.30 元，}相对历史基线提高约 2179 元，最差 10\% 经营利润均值由负转正，收益与尾部风险同时改善。花叶类和辣椒类是主要利润来源，与水生根茎类共同构成补货主体。

\end{enumerate}

\subsection{问题三的模型建立与求解}
\subsubsection{单品需求分解}

对日粒度单品销售表进行日期筛选，仅保留 2023 年 6 月 24 日至 6 月 30 日且毛销售量大于 0 的记录，按单品编码去重后得到候选集合 $I$，共 49 个候选单品，覆盖六个品类。损耗率采用近期销售量加权值，参考价格采用该周销售收入与销量之比；最新批发价取候选周内该单品最后一个非空批发价。

为了与第二问成本预测保持一致，先计算候选周各品类的销量加权批发成本，再以第二问 7 月 1 日品类成本预测相对该近期成本的比例，调整各单品最新批发价：
\begin{equation}
    \hat{w}_i = w_{i,\mathrm{latest}} \times \frac{\hat{w}_{c(i),\mathrm{Q2}}}{\bar{w}_{c(i),\mathrm{recent}}}
\end{equation}
该比例缩放既保留了同品类单品之间的成本差异，又使单品成本整体变化方向与第二问一致。场景成本进一步叠加第二问成本残差：
\begin{equation}
    w_i^{\omega} = \hat{w}_i \times \exp\bigl(e_{c(i)}^{\omega}\bigr)
\end{equation}

对品类 $c$ 内的单品 $i$，记最近 7 日销量为 $Q_i^{\mathrm{recent}}$，品类最近 7 日总销量为 $Q_c^{\mathrm{recent}}$，历史可比星期六长期份额为 $r_i^{\mathrm{long}}$。构造经验贝叶斯份额：
\begin{equation}
    \hat{r}_i = \frac{Q_i^{\mathrm{recent}} + \kappa_c r_i^{\mathrm{long}}}{Q_c^{\mathrm{recent}} + \kappa_c},\qquad \sum_{i:c(i)=c}\hat{r}_i = 1
\end{equation}
式（3）等价于在近期观测上加入 $\kappa_c$ kg 的长期先验''伪样本''。当近期样本充分时，$Q_c^{\mathrm{recent}}$ 远大于 $\kappa_c$，份额主要由最近一周决定；当近期样本有限时，估计值向长期可比日结构收缩。$\kappa_c$ 不主观指定，而是在目标日前的历史星期六上滚动验证，从候选网格中选择预测误差较小的值。

为了形成单品份额场景，从过去三年的可比星期六中按整日有放回抽样，计算抽样日内份额 $r_i^{\mathrm{obs},\omega}$，再与后验份额混合：
\begin{equation}
    r_i^{\omega} = h \cdot r_i^{\mathrm{obs},\omega} + (1-h) \cdot \hat{r}_i,\qquad h = 0.50
\end{equation}
每个场景、每个品类内部重新归一化，确保份额之和严格为 1。保留了同一日各单品份额间的相关变化。

连续定价会使价格响应项产生非线性，也可能给出历史中从未出现过的难执行价格。因此对每个单品建立有限价格网格：优先使用该单品近一年实际加价率的 10\%、30\%、50\%、70\%、90\% 销量加权分位点；若有效销售天数少于 30 天或不同价格少于 5 个，则回退到所属品类历史加价率。所有加价率均截断在第二问估计的品类合理边界内，最终价格按 0.1 元/kg 取整。

对单品 $i$ 的价格方案 $k$，复用第二问所属品类价格弹性 $\varepsilon_c$，构造场景需求：
\begin{equation}
    d_{ik}^{\omega} = r_i^{\omega} D_{c(i)}^{\omega} \left(\frac{p_{ik}}{p_i^{\mathrm{ref}}}\right)^{\varepsilon_{c(i)}}
\end{equation}
由于蔬菜需求通常随价格上升而下降，估计弹性为负时，价格高于参考价会压低需求。式（5）把第二问的品类需求、第三问的单品结构和定价决策连接为统一口径。

设置二元变量 $x_i$ 表示单品是否入选，$y_{ik}$ 表示是否采用价格 $k$；连续变量 $q_i$ 表示订货量，$s_{ik}^{\omega}$ 表示场景销量，$u_c^{\omega}$ 表示品类未满足需求。

商品数量约束为：
\begin{equation}
    27 \leq \sum_{i\in I} x_i \leq 33
\end{equation}
实际求解中分别固定 $K=27,\ldots,33$ 枚举，以便获得完整的组合规模前沿。每个单品至多选择一个价格，未入选时不定价：
\begin{equation}
    \sum_{k\in K_i} y_{ik} = x_i
\end{equation}
最低陈列量与数据尺度化订货上界为：
\begin{equation}
    2.5\,x_i \leq q_i \leq M_i x_i
\end{equation}
其中 $M_i$ 取各价格方案场景需求 99\% 分位数的 1.2 倍，并除以可售比例 $1-\lambda_i$，且不低于 2.5 kg。

第三问是第二问品类计划的单品落实，因此设置总订货量上限：
\begin{equation}
    \sum_{i\in I} q_i \leq Q^{\mathrm{Q2}} = 362.2627
\end{equation}
场景销量同时受需求和损耗后库存限制：
\begin{equation}
    0 \leq s_{ik}^{\omega} \leq d_{ik}^{\omega} y_{ik}
\end{equation}
\begin{equation}
    \sum_{k\in K_i} s_{ik}^{\omega} \leq (1-\lambda_i) q_i
\end{equation}
品类未满足需求满足：
\begin{equation}
    u_c^{\omega} + \sum_{i:c(i)=c}\sum_{k\in K_i} s_{ik}^{\omega} \geq D_c^{\omega},\qquad u_c^{\omega} \geq 0
\end{equation}




\subsubsection{选品、定价和补货联合优化}

若直接最小化总缺货量，大需求的花叶类和辣椒类会支配目标，小体量品类可能被忽视。为体现题目''各品类需求''，定义六品类等权的平均标准化缺货率：
\begin{equation}
    L = \frac{\frac{1}{|\Omega||C|}\sum_{\omega\in\Omega}\sum_{c\in C} u_c^{\omega}}{\max\bigl(D_c^{\omega},\,10^{-9}\bigr)}
\end{equation}
第一阶段求解：
\begin{equation}
    L^* = \min L
\end{equation}
其中 $L$ 为六品类等权平均标准化缺货率，满足以下约束：商品数量约束 $27\leq\sum_i x_i\leq 33$，选品互斥约束 $\sum_k y_{ik}=x_i$，最低陈列量与补货上界 $2.5x_i\leq q_i\leq M_i x_i$，总订货量上限 $\sum_i q_i\leq Q^{\mathrm{Q2}}$，场景销量与需求匹配 $0\leq s_{ik}^\omega\leq d_{ik}^\omega y_{ik}$，场景销量与库存匹配 $\sum_k s_{ik}^\omega\leq (1-\lambda_i)q_i$，以及品类需求满足约束 $u_c^\omega+\sum_i\sum_k s_{ik}^\omega\geq D_c^\omega$。

场景利润为销售收入减去订货成本，成本按损耗前订货量计：
\begin{equation}
    \Pi^{\omega} = \sum_i\sum_k p_{ik} s_{ik}^{\omega} - \sum_i w_i^{\omega} q_i
\end{equation}
第一阶段使用有限代表场景，若强制第二阶段严格等于 $L^*$，可能因微小场景噪声排除大量利润更优方案。因此计算第一阶段各场景服务损失的标准误 $\delta$，并增加约束：
\begin{equation}
    L \leq L^* + \delta
\end{equation}
为避免只优化均值掩盖极端亏损，引入最差 $\alpha=10\%$ 场景平均利润 $\mathrm{LTM}_{\alpha}$。利用辅助变量 $\eta$ 和 $\xi_{\omega}$ 进行线性化：
\begin{equation}
    \xi_{\omega} \geq \eta - \Pi^{\omega},\qquad \xi_{\omega} \geq 0
\end{equation}
\begin{equation}
    \mathrm{LTM}_{\alpha}(\Pi) = \eta - \frac{1}{\alpha|\Omega|}\sum_{\omega\in\Omega}\xi_{\omega}
\end{equation}
第二阶段目标为：
\begin{equation}
    \max\bigl\{(1-\gamma)\,{E}[\Pi] + \gamma\,\mathrm{LTM}_{0.10}(\Pi)\bigr\},\qquad \gamma = 0.25
\end{equation}
并满足约束（6）-（12）及（16）-（18）。该目标不是把缺货和利润任意相加，而是在服务水平已基本锁定后处理风险-收益权衡，符合题目优先级。


表 \ref{tab:q3_frontier} 显示 $K$ 从 27 增至 33 时，完整场景需求满足率由 76.93\% 单调升至 77.81\%，期望利润由 62.52 元增至 163.75 元，风险调整利润由 $-81.57$ 元改善为 17.11 元。在给定的 Q2 补货总量约束下，增加单品并没有挤压利润，反而提高了需求覆盖和价格匹配能力，因此最终选取上限 $K=33$。

\begin{table}[H]
\centering
\begin{tabularx}{\linewidth}{>{\centering\arraybackslash}X>{\centering\arraybackslash}X>{\centering\arraybackslash}X>{\centering\arraybackslash}X>{\centering\arraybackslash}X>{\centering\arraybackslash}X}
    \toprule
    $K$ & 需求满足率/\% & 期望利润/元 & 最差10\%平均利润/元 & 风险调整利润/元 & 总订货/kg \\
    \midrule
    27 & 76.93 & 62.52 & $-513.82$ & $-81.57$ & 362.263 \\
    28 & 77.00 & 95.04 & $-487.70$ & $-50.64$ & 362.263 \\
    29 & 77.06 & 110.48 & $-475.46$ & $-36.01$ & 362.263 \\
    30 & 77.37 & 126.92 & $-461.86$ & $-20.27$ & 362.263 \\
    31 & 77.53 & 138.85 & $-448.61$ & $-8.01$ & 362.263 \\
    32 & 77.67 & 149.85 & $-435.68$ & 3.47 & 362.263 \\
    \textbf{33} & \textbf{77.81} & \textbf{163.75} & \textbf{$-422.81$} & \textbf{17.11} & \textbf{362.263} \\
    \bottomrule
\end{tabularx}
\caption{组合规模前沿}
\label{tab:q3_frontier}
\end{table}

\begin{figure}[H]
\centering
\includegraphics[width=\textwidth]{figures/fig_q3_k_frontier.png}
\caption{经营单品数量的服务---收益前沿图}
\label{fig:q3_frontier}
\end{figure}


\textbf{最终方案选择 33 个单品，总订货量 362.2627 kg}。完整 600 场景评价结果：\textbf{需求满足率 77.81\%}，期望利润 \textbf{163.75 元}，最差 10\% 平均利润 $-422.81$ 元，风险调整利润 17.11 元。

透明基线（近期销量前 33 名 + 中位价格 + 近期日均订货）满足率 64.61\%，期望利润 56.48 元。\textbf{主模型相对基线满足率提高 13.20 个百分点，期望利润增加 107.27 元}，改进来自品类覆盖、单品结构和定价的联合优化，而非单纯扩大订货量。


\begin{table}[H]
\centering
\begin{tabularx}{\linewidth}{>{\centering\arraybackslash}X>{\centering\arraybackslash}X>{\centering\arraybackslash}X>{\centering\arraybackslash}X>{\centering\arraybackslash}X>{\centering\arraybackslash}X>{\centering\arraybackslash}X>{\centering\arraybackslash}X}
    \toprule
    品类 & 候选数 & 入选数 & 期望需求/kg & 期望销量/kg & 满足率/\% & 订货量/kg & 期望利润/元 \\
    \midrule
    水生根茎类 & 7 & 4 & 17.63 & 14.91 & 91.78 & 25.87 & $-66.31$ \\
    花叶类 & 17 & 11 & 195.47 & 85.56 & 55.12 & 116.03 & 158.98 \\
    花菜类 & 2 & 2 & 16.90 & 14.94 & 94.46 & 23.68 & $-32.22$ \\
    茄类 & 5 & 4 & 33.50 & 26.71 & 90.86 & 45.15 & $-9.40$ \\
    辣椒类 & 10 & 6 & 135.22 & 63.55 & 66.26 & 90.94 & 119.87 \\
    食用菌 & 8 & 6 & 71.09 & 37.66 & 68.37 & 60.60 & $-7.16$ \\
    \bottomrule
\end{tabularx}
\caption{品类层面结果}
\label{tab:q3_cat}
\end{table}

由表 \ref{tab:q3_cat}：水生根茎、花菜和茄类的满足率均超过 90\%，说明较小品类没有被大品类目标吞没；花叶类和辣椒类需求体量大，在总订货上限下仍是主要缺口来源。部分品类期望利润为负是服务优先目标的直接结果：模型为维持品类覆盖保留了若干低毛利或高损耗单品。整体利润仍为正，且显著高于基线。

\begin{figure}[H]
\centering
\includegraphics[width=\textwidth]{figures/fig_q3_category_strategy.png}
\caption{六品类选品、补货、服务与利润概览}
\label{fig:q3_overview}
\end{figure}

由图~\ref{fig:q3_overview}说明：

\begin{enumerate}[leftmargin=2em,itemsep=2pt,parsep=0pt]

    \item \textbf{服务与利润同步提升。}
    相比基线，需求满足率提高 $13.20$ 个百分点，期望利润增加 $107.27$ 元，说明联合优化并非简单增加订货或提高价格。

    \item \textbf{品类间存在资源竞争。}
    花叶类需求占比最高但满足率仅为 $55.12\%$，模型采用六品类等权缺货率进行优化，以保证各品类获得基本需求覆盖。

    \item \textbf{尾部风险仍需关注。}
    虽然期望利润为正，但最差 $10\%$ 场景平均利润为 $-422.81$ 元，说明极端成本与需求波动下仍存在亏损风险。

    \item \textbf{单品需求分配具有较大影响。}
    历史份额权重变化会显著影响满足率和利润，表明单品需求分解是模型的重要环节，实际经营中应进一步采集库存、缺货和陈列等数据。

\end{enumerate}


\subsection{模型检验与灵敏度}

\textbf{可行性与一致性检验。} 独立输出校验器重新从原始流水识别候选集合，并检查候选唯一性、成本与损耗范围、品类份额和为 1、价格网格来源、$K=27-33$ 前沿、最低订货量、总量上限、销量不超过需求、未售量恒等式、六品类汇总与单品汇总勾稽、利润汇总和敏感性参数范围。共 77 项检查全部通过。另有 17 项单元测试覆盖日期边界、经验贝叶斯收缩、价格回退、损耗口径、MILP 约束、第一阶段复用和缓存失效拒绝，全部通过。

\textbf{代表场景泛化检验。} 优化使用 60 个代表场景，最终在 600 个完整场景上评价。代表场景的第二阶段服务损失为 0.214430，完整场景服务损失为 0.221921，绝对差为 0.007490，即 0.749 个百分点，明显低于预设阈值 0.03。这表明场景压缩没有造成显著的服务水平高估。

\textbf{六折场景稳定性检验。} 将 600 个场景分为 6 折，每折 100 个场景，分别计算固定最终策略的满足率和利润。六折满足率依次为 77.26\%、76.72\%、78.91\%、77.52\%、79.02\% 和 77.41\%，均值 77.81\%，样本标准差 0.009405，低于 0.05 的稳定性阈值；最小与最大满足率仅相差约 2.31 个百分点。六折期望利润标准差为 31.69 元，反映利润对成本与需求联合波动比服务率更敏感，因此模型报告尾部利润是必要的。

\textbf{基线有效性检验。} 与近期销量排序基线相比，主模型的满足率、期望利润和最差 10\% 平均利润三个核心指标均改善。由于基线规则透明、使用相同候选、相同价格可行域和同一 600 场景评价，比较不依赖隐藏的外部数据，能够证明联合优化相对常规经验策略具有实际增益。

灵敏度分析采用''改变一个参数并完整重求解''的方式，参数包括风险权重 $\gamma$、价格弹性缩放、历史份额权重 $h$ 和补货上限系数。表 \ref{tab:q3_sens} 汇总 9 个不重复方案（见附录）。

所有 9 个方案的最差 10\% 平均利润均为负，范围为 $-599.80$ 至 $-264.99$ 元。因此灵敏度分析的合格结论应表述为：\textbf{主方案在服务率与选品结构上具有较好稳定性，且明显优于朴素基线；但单品份额结构和极端场景亏损仍是主要风险，实际应用应结合当天到货价与开店后早期销量进行滚动修正。}

\begin{figure}[H]
\centering
\includegraphics[width=\textwidth]{figures/fig_q3_model_comparison.png}
\caption{主模型与透明基线比较图}
\label{fig:q3_baseline}
\end{figure}





第二问提供 600 个保留跨品类相关性的联合场景。为控制含场景销量变量的 MILP 规模，按六品类总需求从低到高排序，在等距分位位置选择 60 个代表场景，且保留两端极端场景。对 $K=27,\ldots,33$ 分别执行两阶段优化，再将每个方案放回完整 600 场景计算满足率、期望利润和最差 10\% 平均利润，最终按服务优先原则选取 $K$。

\begin{figure}[H]
\centering
\includegraphics[width=\textwidth]{figures/fig_q3_sensitivity.png}
\caption{关键参数灵敏度与选品稳定性图}
\label{fig:q3_sensitivity}
\end{figure}



通过采用经验贝叶斯动态收缩估计单品需求份额，以历史可比星期六份额为长程先验、近期销量为后验信息，对 $K=27$ 至 $33$ 执行两阶段词典序优化并在 600 个完整场景上评价。\textbf{最终选取 $K=33$，共 33 个可售单品、总订货 362.26 kg，需求满足率 77.81\%，期望利润 163.75 元，}较近期销量基线满足率提高 13.20 个百分点、利润增加 107.27 元，显著优于常规经验策略。花叶类获得最大补货份额，多数单品定价位于参考价附近，少量高定价单品通过牺牲销量换取单品利润贡献。得到\textbf{结果如图 \ref{fig:q3_resource} 和图 \ref{fig:q3_sku} 所示。}
\begin{figure}[H]
\centering
\includegraphics[width=0.8\textwidth]{figures/fig_q3_category_radial_allocation.png}
\caption{六品类需求覆盖与补货资源分配图}
\label{fig:q3_resource}
\end{figure}

\begin{figure}[H]
\centering
\includegraphics[width=0.8\textwidth]{figures/fig_q3_sku_strategy_bubbles.png}
\caption{单品订货---售价---销量组合图}
\label{fig:q3_sku}
\end{figure}
\subsection{问题四的求解}
针对问题四，本文建议商超重点采集以下几类数据以提升补货与定价决策的准确性：



\textbf{（1）门店客流数据。}多品类销量同涨同跌可能受客流总量变化驱动，而非品类间的真实替代关系。建议采集每日门店客流量，可在问题一的相关关系分析中排除共同因素干扰，提高品类间相关关系判断的准确性。

\textbf{（2）天气与节假日信息。}蔬菜销量受天气和节假日影响显著，建议补充历史同期气象数据（温度、降雨）和节假日标记，用于改进问题二中需求预测的季节性调整和问题三中单品份额的动态估计。

\subsection{模型的优点}
本模型的主要优点包括以下几个方面：

（1）\textbf{建模框架系统性强。}
 四个问题层层衔接，形成从分布分析、需求预测到单品优化和数据改进的完整建模链路。

（2）\textbf{需求预测方法贴合实际。}
预测方法贴合实际。 结合 GBDT、季节朴素模型与 DML，兼顾需求预测的季节性与价格响应关系。

（3）\textbf{场景生成兼顾相关性与多样性。}
场景生成合理。 通过联合残差保留品类间相关性，使生成的需求场景更符合实际销售波动特征。

（4）\textbf{优化目标兼顾收益与风险。}
引入均值—下尾 CVaR，在提高期望收益的同时降低极端亏损风险。

\subsection{模型的缺点}


本模型存在以下局限性和可改进之处：

（1）\textbf{问题一的分布模型仅考虑了四种常用分布。}虽然在所选分布中通过 AIC 和 KS 检验进行了优化，但实际销量分布可能包含更复杂的结构（如混合分布、阈值效应），四种候选分布未必能覆盖所有品类的真实分布形态。

（2）\textbf{价格弹性估计对样本量敏感。}部分品类在历史数据中的价格变动幅度有限，DML 估计的斜率可能存在较大的标准误，导致条件价格响应的置信区间较宽，进而影响问题二中加成率设定的精度。
%参考文献

\begin{thebibliography}{9}

\bibitem{buzby2015produce}
Buzby J C, Bentley J T, Padera B, et al.
Estimated fresh produce shrink and food loss in U.S. supermarkets[J].
\textit{Agriculture},
2015, 5(3): 626--648.

\bibitem{broekmeulen2019foodwaste}
Broekmeulen R A C M, van Donselaar K H.
Quantifying the potential to improve on food waste, freshness and sales for perishables in supermarkets[J].
\textit{International Journal of Production Economics},
2019, 209: 265--273.

\bibitem{akaike1974}
Akaike H.
A new look at the statistical model identification[J].
\textit{IEEE Transactions on Automatic Control},
1974, 19(6): 716--723.


\bibitem{marshall1920}
Marshall A.
\textit{Principles of Economics: An Introductory Volume}[M].
8th ed.
London: Macmillan, 1920.

\bibitem{tellis1988}
Tellis G J.
The Price Elasticity of Selective Demand: A Meta-Analysis of Econometric Models of Sales[J].
\textit{Journal of Marketing Research}, 1988, 25(4): 331--341.
DOI: 10.1177/002224378802500401.


\end{thebibliography}

\newpage
%附录
\begin{appendices}




% 附录 A: 问题二品类逐日补货与定价策略
\section{问题二品类逐日补货与定价策略}
\label{app:q2_strategy}


\begin{longtable}{*{7}{>{\centering\arraybackslash}p{\dimexpr\textwidth/7-2\tabcolsep\relax}}}
\caption{2023 年 7 月 1—7 日各品类日补货总量与定价策略}\label{tab:q2_strategy}\\
\toprule
日期 & 品类 & 成本/元 & 加成率 & 售价/元 & 补货/kg & 利润/元 \\
\midrule
\endfirsthead
\multicolumn{7}{c}{\tablename~\thetable{} -- 续表} \\
\toprule
日期 & 品类 & 成本/元 & 加成率 & 售价/元 & 补货/kg & 利润/元 \\
\midrule
\endhead
\endfoot
\bottomrule
\endlastfoot
07-01 & 水生根茎 & 11.64 & 64.2\% & 19.11 & 12.85 & 35.67 \\
07-01 & 花叶类 & 3.41 & 88.8\% & 6.44 & 156.64 & 305.64 \\
07-01 & 花菜类 & 8.67 & 74.6\% & 15.14 & 13.69 & 47.51 \\
07-01 & 茄类 & 4.66 & 86.5\% & 8.69 & 27.86 & 63.95 \\
07-01 & 辣椒类 & 3.36 & 101.1\% & 6.76 & 101.01 & 209.72 \\
07-01 & 食用菌 & 3.65 & 69.2\% & 6.18 & 50.21 & 88.06 \\
07-02 & 水生根茎 & 12.69 & 65.5\% & 21.00 & 9.41 & 30.60 \\
07-02 & 花叶类 & 3.43 & 100.4\% & 6.87 & 121.94 & 267.88 \\
07-02 & 花菜类 & 8.10 & 89.0\% & 15.31 & 13.60 & 54.32 \\
07-02 & 茄类 & 4.28 & 92.2\% & 8.23 & 23.33 & 53.05 \\
07-02 & 辣椒类 & 3.15 & 127.7\% & 7.18 & 76.67 & 193.65 \\
07-02 & 食用菌 & 3.69 & 70.9\% & 6.30 & 36.09 & 67.32 \\
07-03 & 水生根茎 & 12.49 & 65.5\% & 20.66 & 7.53 & 23.26 \\
07-03 & 花叶类 & 3.46 & 100.4\% & 6.94 & 91.34 & 204.72 \\
07-03 & 花菜类 & 7.89 & 81.5\% & 14.32 & 10.82 & 39.58 \\
07-03 & 茄类 & 4.09 & 92.2\% & 7.85 & 19.32 & 43.30 \\
07-03 & 辣椒类 & 3.23 & 132.6\% & 7.50 & 75.75 & 209.62 \\
07-03 & 食用菌 & 3.47 & 70.9\% & 5.93 & 37.20 & 64.97 \\
07-04 & 水生根茎 & 12.84 & 65.5\% & 21.25 & 8.06 & 25.36 \\
07-04 & 花叶类 & 3.42 & 100.4\% & 6.86 & 124.92 & 278.30 \\
07-04 & 花菜类 & 7.77 & 87.1\% & 14.54 & 13.44 & 53.26 \\
07-04 & 茄类 & 4.71 & 75.9\% & 8.28 & 7.23 & 15.48 \\
07-04 & 辣椒类 & 3.47 & 132.6\% & 8.06 & 69.89 & 204.32 \\
07-04 & 食用菌 & 3.46 & 70.9\% & 5.91 & 35.06 & 59.92 \\
07-05 & 水生根茎 & 11.93 & 65.5\% & 19.74 & 7.14 & 21.69 \\
07-05 & 花叶类 & 3.49 & 100.4\% & 7.00 & 127.32 & 285.47 \\
07-05 & 花菜类 & 7.83 & 88.0\% & 14.72 & 13.99 & 53.06 \\
07-05 & 茄类 & 4.73 & 89.6\% & 8.96 & 14.42 & 35.81 \\
07-05 & 辣椒类 & 3.40 & 132.6\% & 7.90 & 70.61 & 196.38 \\
07-05 & 食用菌 & 3.21 & 70.9\% & 5.49 & 48.70 & 74.76 \\
07-06 & 水生根茎 & 12.00 & 65.5\% & 19.85 & 11.21 & 34.73 \\
07-06 & 花叶类 & 3.35 & 100.4\% & 6.72 & 138.03 & 289.61 \\
07-06 & 花菜类 & 8.04 & 91.9\% & 15.44 & 19.39 & 81.99 \\
07-06 & 茄类 & 4.52 & 81.6\% & 8.21 & 10.57 & 23.64 \\
07-06 & 辣椒类 & 3.64 & 132.6\% & 8.47 & 86.38 & 266.91 \\
07-06 & 食用菌 & 3.22 & 70.9\% & 5.50 & 44.68 & 70.34 \\
07-07 & 水生根茎 & 11.69 & 65.5\% & 19.35 & 9.70 & 26.15 \\
07-07 & 花叶类 & 3.16 & 100.4\% & 6.33 & 139.85 & 276.71 \\
07-07 & 花菜类 & 7.87 & 91.9\% & 15.10 & 21.70 & 91.02 \\
07-07 & 茄类 & 4.66 & 92.2\% & 8.96 & 21.19 & 57.58 \\
07-07 & 辣椒类 & 3.72 & 132.6\% & 8.65 & 79.45 & 248.22 \\
07-07 & 食用菌 & 4.00 & 70.9\% & 6.84 & 40.13 & 76.77 \\
\end{longtable}

\section{六品类月度季节指数曲线}
\begin{figure}[H]
    \centering
    \includegraphics[width=1.2\textwidth]{figures/Q1蔬菜品类月度季节指数.png}
    \caption{ 六品类月度季节指数曲线}
    \label{fig:Person}
\end{figure}
\section{关键参数灵敏度分析}

\begin{table}[H]
\centering
\begin{tabularx}{\linewidth}{>{\centering\arraybackslash}X>{\centering\arraybackslash}X>{\centering\arraybackslash}X>{\centering\arraybackslash}X>{\centering\arraybackslash}X>{\centering\arraybackslash}X}
    \toprule
    方案 & 参数变化 & 满足率/\% & 期望利润/元 & 最差10\%利润/元 & 选品Jaccard \\
    \midrule
    基准 & $\gamma$=0.25，$h$=0.50，上限=1.00 & 77.81 & 163.75 & $-422.81$ & 1.000 \\
    风险权重 0 & $\gamma$=0 & 77.82 & 167.53 & $-432.34$ & 1.000 \\
    风险权重 0.5 & $\gamma$=0.50 & 78.01 & 151.45 & $-421.63$ & 1.000 \\
    弹性 $\times$0.8 & $\varepsilon\times$0.8 & 77.72 & 178.46 & $-411.04$ & 1.000 \\
    弹性 $\times$1.2 & $\varepsilon\times$1.2 & 77.94 & 153.52 & $-432.57$ & 1.000 \\
    历史份额权重 0.25 & $h$=0.25 & 82.83 & 356.56 & $-268.75$ & 0.941 \\
    历史份额权重 0.75 & $h$=0.75 & 71.86 & $-51.83$ & $-599.80$ & 0.833 \\
    补货上限 0.85 & $Q$上限$\times$0.85 & 74.44 & 235.22 & $-264.99$ & 1.000 \\
    补货上限 0.925 & $Q$上限$\times$0.925 & 76.27 & 197.42 & $-345.23$ & 1.000 \\
    \bottomrule
\end{tabularx}
\caption{关键参数灵敏度分析}
\label{tab:q3_sens}
\end{table}
\end{appendices}
\end{document}
